\section{PCD_to_polar}


\begin{frame}[fragile]{Implementation of the polar coordinates decoupling (PCD). To cartesian}
    \begin{columns}[T,onlytextwidth]
        \column{0.55\textwidth}
        \begin{adjustbox}{max height=0.65\textheight,keepaspectratio}
\begin{lstlisting}[language=Python, basicstyle=\ttfamily\scriptsize, aboveskip=2pt, belowskip=2pt, numbers=left, numberstyle=\tiny, frame=single, breaklines=true]
def get_unit_directions(directions: torch.Tensor) -> torch.Tensor:
    sin_matrix = torch.sin(directions)
    cos_matrix = torch.cos(directions)
    cum_sin = torch.cumprod(sin_matrix, dim=-1)
    sin_prefix = torch.cat([torch.ones_like(sin_matrix[:, :1]), cum_sin[:, :-1]], dim=-1)
    x_except_last = sin_prefix * cos_matrix
    x_last = cum_sin[:, -1:].clone()
    unitary_directions = torch.cat([x_except_last, x_last], dim=-1)
    return unitary_directions

def to_cartesian(phis: torch.Tensor, r: torch.Tensor) -> torch.Tensor:
    if r.ndim == 1:
        r = r.unsqueeze(-1)
    unit_dirs = PCDVQ.get_unit_directions(phis)
    x = r * unit_dirs
    return x
\end{lstlisting}
        \end{adjustbox}
    \column{0.4\textwidth}
        $get\_unit\_directions$: constructs a unit direction vector based on sets of angles in hyperspherical coordinates. Takes $\cos$ for each axis, accumulates the products of $\sin $ from left to right, and obtains the components as $sin\_prefix\cdot\cos $, and the last one as the full product of $\sin $\\
        $to\_cartesian$: сonverts the radius $r$ to a column if necessary, calculates unit directions via $get\_unit\_directions(phis)$, and multiplies them by $r$\\
  \end{columns}
\end{frame}
